%! Random Sampling Methods — Unit 3, Lesson 3.2 — Slides (Beamer)
\documentclass[aspectratio=169, 11pt]{beamer}
\usepackage{apstats-beamer}   % provides \navyheader{} and \sectionlabel[color]{}

\begin{document}

% TITLE SLIDE (CourseName is not defined in beamer, so the name is literal)
{
\setbeamercolor{background canvas}{bg=navy}
\begin{frame}[plain]
  \vspace{2.8cm}\hspace{0.55cm}%
  \begin{minipage}{0.9\textwidth}
    {\color{goldacc}\bfseries\small LESSON 3.2}\\[0.5cm]
    {\color{white}\bfseries\LARGE Random Sampling Methods}\\[1.2cm]
    {\color{skymid}\small Statistics~$\cdot$~Unit 3}
  \end{minipage}
\end{frame}
}

% GRADUAL-RELEASE deck: targets → warm-up → I DO divider → one frame per instruction row of the
% notes (SRS; Strata & Clusters; Choosing, which carries the crux) → WE DO (Parking at
% Riverside) → spoken debrief → close & START THE HOMEWORK, which is the individual practice.
% There is no group activity and no independent practice set in the notes.
\newcommand{\redink}[1]{{\color{redacc}\bfseries #1}}

% ── LEARNING TARGETS ─────────────────────────────────────────────────────────
\begin{frame}[t, plain]
  \navyheader{Today's targets}
  \sectionlabel{Lesson 3.2}
  By the end of today, I can\ldots
  \begin{itemize}\setlength{\itemsep}{5pt}
    \item describe how to choose a \textbf{simple random sample}, and tell sampling
          \textbf{with replacement} from \textbf{without replacement}.
    \item identify a \textbf{stratified random sample} and a \textbf{cluster sample}.
    \item explain, in context, why one sampling method is a better choice than another.
  \end{itemize}
  \begin{block}{How today works}
    Warm-up $\to$ \textbf{notes together} $\to$ \textbf{one we work as a class} $\to$
    debrief $\to$ \textbf{you start the homework here, alone}.
  \end{block}
\end{frame}

% ── WARM-UP ──────────────────────────────────────────────────────────────────
\begin{frame}[t, plain]
  \navyheader{Warm-up}
  \sectionlabel{5 minutes $\cdot$ on your own}
  Lakeview High: 1{,}200 students, 300 per grade. The principal will survey 80 about hours
  worked at a paid job.
  \begin{enumerate}\setlength{\itemsep}{6pt}
    \item Population? Sample?
    \item All 1{,}200 names in a hat, draw 80. Same chance for everyone?
    \item The hat could give very few seniors --- and seniors work more. So? How could you make
          sure every grade is in?
    \item Why might surveying 2 whole advisory groups be \emph{easier}?
  \end{enumerate}
  \begin{block}{Keep items 3 and 4 in mind}
    Item 3 is a \emph{stratified} sample. Item 4 is a \emph{cluster} sample. The notes name
    them.
  \end{block}
\end{frame}

% ── I DO — direct instruction begins ─────────────────────────────────────────
{
\setbeamercolor{background canvas}{bg=navy}
\begin{frame}[plain]
  \vspace{1.8cm}\hspace{0.8cm}%
  \begin{minipage}{0.86\textwidth}
    {\color{goldacc}\bfseries\small GUIDED NOTES \ \textbullet\ \ I DO}\\[0.7cm]
    {\color{white}\bfseries\Large Open to your Guided Notes. Fill each blank as we name it
    --- not before.}\\[0.9cm]
    {\color{skymid}\small Four terms go in the vocabulary box today: simple random sample,
    without replacement, stratified random sample, cluster sample.}
  \end{minipage}
\end{frame}
}

% ── ROW 1: SRS ───────────────────────────────────────────────────────────────
\begin{frame}[t, plain]
  \navyheader{The simple random sample}
  \sectionlabel{notes, SRS $\cdot$ problems 1--3}
  \vspace{2pt}
  \begin{columns}[t]
    \begin{column}{0.48\textwidth}
      \begin{block}{Simple random sample (SRS)}
        Every \redink{group} of $n$ individuals has an \redink{equal} chance of being the
        sample.
      \end{block}
      \vspace{3pt}
      \begin{block}{Without / with replacement}
        Without: each individual can be chosen only \redink{once}. With: it can be chosen
        again.
      \end{block}
    \end{column}
    \begin{column}{0.48\textwidth}
      {\small\textbf{How to choose one}}
      \begin{enumerate}\setlength{\itemsep}{3pt}\small
        \item \textbf{Label} each individual 1 to $N$.
        \item \textbf{Randomize} with a random number generator or a random digit table.
        \item \textbf{Select}: skip numbers out of range and \redink{repeats} until you have
              $n$.
      \end{enumerate}
      \vspace{6pt}
      {\small Digits, four at a time, students 0001--1200:\\
      \texttt{0738 1195 9661 0738 0452}}
    \end{column}
  \end{columns}
\end{frame}

% ── ROW 2: STRATA & CLUSTERS ─────────────────────────────────────────────────
\begin{frame}[t, plain]
  \navyheader{Strata and clusters}
  \sectionlabel{notes, Strata \& Clusters $\cdot$ problems 4--6}
  \vspace{2pt}
  \begin{columns}[t]
    \begin{column}{0.48\textwidth}
      \begin{block}{Stratified random sample}
        Split into \textbf{strata}, groups \redink{alike} inside. Take an SRS from
        \redink{every} stratum; combine.\\[2pt]
        {\small Lakeview: 20 from each grade.}
      \end{block}
    \end{column}
    \begin{column}{0.48\textwidth}
      \begin{block}{Cluster sample}
        Split into \textbf{clusters}, groups that each look like the \redink{population}.
        Randomly choose some; measure \redink{everyone} in them.\\[2pt]
        {\small Lakeview: 2 of the 48 mixed advisory groups.}
      \end{block}
    \end{column}
  \end{columns}
  \vspace{8pt}
  {\small Named only: a \textbf{systematic random sample} takes a random start, then every
  $k$th individual on a list.}
  \vspace{6pt}
  \begin{block}{Problem 6}
    Pick 1 grade at random, survey all 300. Each student has a 1-in-4 chance. Is it an SRS?
  \end{block}
\end{frame}

% ── ROW 3: CHOOSING — THE CRUX ───────────────────────────────────────────────
\begin{frame}[t, plain]
  \navyheader{Choosing a method}
  \sectionlabel{notes, Choosing $\cdot$ problems 7--10 $\cdot$ the whole point of today}
  \vspace{4pt}
  \begin{columns}[t]
    \begin{column}{0.46\textwidth}
      \begin{block}{Stratify when\ldots}
        the response \redink{differs} between groups. Alike inside $\Rightarrow$ the estimate
        varies less than an SRS.
      \end{block}
    \end{column}
    \begin{column}{0.46\textwidth}
      \begin{block}{Cluster when\ldots}
        reaching scattered people is \redink{expensive}. Cheaper --- best when each cluster
        looks like the whole.
      \end{block}
    \end{column}
  \end{columns}
  \vspace{8pt}
  Jalen: ``Use the grades as clusters --- randomly pick 2 grades and survey everyone in
  them.'' \emph{What is wrong?}
  \vspace{8pt}
  \begin{block}{Say it, write it in the box}
    Strata are alike inside: sample from \redink{every} one. Clusters each look like the whole:
    take \redink{all} of \redink{some} of them.
  \end{block}
\end{frame}

% ── WE DO ────────────────────────────────────────────────────────────────────
\begin{frame}[t, plain]
  \navyheader{Guided Practice: Parking at Riverside}
  \sectionlabel[goldacc]{We do $\cdot$ you hold the pen $\cdot$ problems 11--14}
  480 apartments in 40 buildings. Every building: 4 one-bedroom, 4 two-bedroom, 4 three-bedroom.
  Mean number of cars per household?
  \begin{itemize}\setlength{\itemsep}{4pt}
    \item Describe an SRS of 60 apartments with a random number generator.
    \item Plan B: 5 random buildings, every apartment. Name it; one advantage.
    \item Plan C: an SRS of 20 of each apartment type. Name it. Why might it beat an SRS?
    \item Plan D: ``The three types are clusters --- pick one type, survey all 160.'' Respond.
  \end{itemize}
\end{frame}

% ── DEBRIEF ──────────────────────────────────────────────────────────────────
\begin{frame}[t, plain]
  \navyheader{Debrief: alike inside, or like the whole?}
  \sectionlabel{8 minutes $\cdot$ pens down, papers up --- we say these aloud}
  \vspace{6pt}
  \begin{columns}[t]
    \begin{column}{0.46\textwidth}
      \begin{block}{Strata --- sample from every one}
        Lakeview's grades.\\
        Riverside's apartment types.
      \end{block}
    \end{column}
    \begin{column}{0.46\textwidth}
      \begin{block}{Clusters --- all of some}
        Lakeview's advisory groups.\\
        Riverside's buildings.
      \end{block}
    \end{column}
  \end{columns}
  \vspace{12pt}
  \begin{block}{Two things to say out loud}
    Why is picking 2 whole grades a bad cluster sample? And why do mixed advisory groups make
    useless strata?
  \end{block}
\end{frame}

% ── CLOSE & START THE HOMEWORK ──────────────────────────────────────────────
\begin{frame}[t, plain]
  \navyheader{Homework --- start it now}
  \sectionlabel[goldacc]{6 items $\cdot$ scored 1 to 5 $\cdot$ due the first class after two study halls}
  \begin{itemize}\setlength{\itemsep}{5pt}
    \item Part A: a hospital's 720 nurses on three shifts and 12 floors --- name each plan, then
          an AP-style multiple choice with a sentence of justification.
    \item Part B: describe a simple random sample, give an advantage of the floor plan, and
          answer a supervisor who wants the shifts as clusters.
  \end{itemize}
  \begin{block}{Silent and alone --- I am circulating}
    \emph{Part B, item 3} is the one I am reading over your shoulder. The AP Practice page is
    \textbf{extra credit, scored 1 to 5} --- work it for the reps.
  \end{block}
\end{frame}

% ── CLOSE ────────────────────────────────────────────────────────────────────
{
\setbeamercolor{background canvas}{bg=navy}
\begin{frame}[plain]
  \vspace{1.5cm}\hspace{0.8cm}%
  \begin{minipage}{0.86\textwidth}
    {\color{goldacc}\bfseries\small BEFORE YOU GO}\\[0.7cm]
    {\color{white}\bfseries\Large Strata are alike inside: sample from every one. Clusters
    each look like the whole: take all of some.}\\[0.9cm]
    {\color{skymid}\small Next: Lesson 3.3 --- Problems with Sampling: who never gets a chance,
    who will not answer, and how wording pushes the answers.}
  \end{minipage}
\end{frame}
}

\end{document}
